\chapter{Sleep Apnoea}
\index{Sleep Apnoea}

\section*{Definition and Overview}
Sleep apnoea is a sleep-related breathing disorder in which breathing repeatedly stops and starts during sleep. The most common form is \textbf{obstructive sleep apnoea} (OSA), in which the muscles of the throat relax and the soft tissues of the upper airway collapse during sleep, intermittently blocking airflow. Less commonly, \textbf{central sleep apnoea} (CSA) occurs when the brain fails to send the correct signals to the muscles that control breathing, and some people have a mixed pattern.

Each pause in breathing, or apnoea, typically lasts from 10 seconds to more than a minute and can occur many times each hour. These episodes lower oxygen levels in the blood and cause repeated brief awakenings that fragment sleep, even though the person is usually unaware of them. As a result, sleep becomes unrefreshing and daytime function is impaired. Sleep apnoea is a chronic condition, but it is highly treatable, and effective management substantially improves symptoms and reduces long-term cardiovascular risk.

\section*{Epidemiology}
Sleep apnoea is one of the most common sleep disorders and is significantly under-diagnosed. It is estimated that around one in seven adults has at least mild obstructive sleep apnoea, with the figure being higher in men, in older adults, and in people who are overweight or obese. Obesity is the strongest risk factor, and rising rates of obesity have been accompanied by increasing prevalence of the condition. It is also more common in people with type 2 diabetes, hypertension, heart failure, and atrial fibrillation. Children can also be affected, most often when enlarged tonsils and adenoids narrow the airway. Because many people are unaware of their night-time breathing pauses, a large proportion of cases remain unrecognised, and the true population burden is probably higher than reported figures suggest.

\section*{Aetiology and Pathophysiology}
During sleep, the muscles that hold the upper airway open, including the muscles of the pharynx and tongue, relax. In obstructive sleep apnoea the airway narrows or collapses completely despite continued efforts to breathe, so airflow is blocked while the chest and abdomen continue to move. Blood oxygen levels fall and carbon dioxide rises, which triggers a brief arousal from sleep that restores muscle tone and reopens the airway. This cycle of collapse, arousal, and recovery repeats many times through the night, producing sympathetic nervous system activation, surges in heart rate and blood pressure, and chronic low-grade inflammation.

\begin{itemize}
    \item \textbf{Obesity:} fat deposition around the neck and pharynx narrows the airway and is the single strongest risk factor
    \item \textbf{Craniofacial anatomy:} a small or receding jaw, a large tongue, a long soft palate, or a crowded upper airway
    \item \textbf{Upper airway narrowing:} large tonsils and adenoids, especially in children, and nasal obstruction
    \item \textbf{Alcohol and sedatives:} enhance relaxation of the pharyngeal muscles and prolong apnoeas
    \item \textbf{Smoking:} causes oedema and inflammation of the upper airway lining
    \item \textbf{Hormonal factors:} men and post-menopausal women are affected more often, suggesting a protective role of female sex hormones
    \item \textbf{Family history:} facial structure and ventilatory control can run in families
    \item \textbf{Central causes:} heart failure, stroke, and opioid use can impair the brainstem's control of breathing
\end{itemize}

\section*{Clinical Features}
The condition is frequently first recognised by a bed partner or family member who observes loud snoring and pauses in breathing. The person themselves often complains mainly of unrefreshing sleep and daytime sleepiness.

\begin{itemize}
    \item Loud, habitual snoring that is often punctuated by gasping, choking, or snorting as breathing resumes
    \item Witnessed pauses in breathing during sleep
    \item Excessive daytime sleepiness, including falling asleep during meetings, at work, or while driving
    \item Waking unrefreshed despite an apparently adequate number of hours in bed
    \item Morning headaches, dry mouth, or a sore throat on waking
    \item Nocturia, or waking frequently overnight to pass urine
    \item Irritability, poor concentration, memory difficulties, and depressed mood
    \item Reduced libido and erectile dysfunction in some men
    \item In children: restless sleep, mouth breathing, snoring, bedwetting, and behavioural or attention problems that may be mistaken for attention deficit
\end{itemize}

\section*{Differential Diagnosis}
Several conditions can produce similar daytime sleepiness or disturbed night-time breathing and should be considered before the diagnosis is made.

\begin{itemize}
    \item Primary (simple) snoring without significant apnoea or oxygen desaturation
    \item Chronic insomnia or inadequate sleep quantity due to lifestyle factors
    \item Narcolepsy, which features uncontrollable sleep attacks and cataplexy
    \item Restless legs syndrome and periodic limb movement disorder
    \item Shift-work disorder and other causes of chronic sleep deprivation
    \item Central sleep apnoea, which requires a separate approach to management
    \item Hypothyroidism, depression, and chronic fatigue syndrome, which can mimic excessive sleepiness
\end{itemize}

\section*{When to Seek Medical Care}
A person should see a doctor if they snore heavily, are told that they stop breathing during sleep, or feel excessively sleepy during the day, particularly if sleepiness affects driving, work, or school. Assessment is also warranted when a child snores persistently, breathes through the mouth at night, or shows behavioural or school problems, since tonsil and adenoid enlargement is often the cause.

\textbf{Contact your local emergency services for:}
\begin{itemize}
    \item Chest pain or severe breathlessness, which may indicate strain on the heart from untreated apnoea
    \item Fainting, collapse, or sudden loss of consciousness
    \item A child with severe difficulty breathing, gasping for air, or turning blue during sleep
    \item Sudden severe headache, weakness on one side of the body, or difficulty speaking, which may signal a stroke
\end{itemize}

\section*{Diagnosis and Investigations}
The diagnosis is based on a characteristic combination of symptoms and is confirmed by an objective sleep study that measures breathing, oxygen levels, and sleep architecture.

\begin{itemize}
    \item \textbf{History:} snoring, witnessed apnoeas, daytime sleepiness, nocturia, morning symptoms, and risk factors including weight gain
    \item \textbf{Epworth Sleepiness Scale:} a validated questionnaire that quantifies the tendency to fall asleep during daily activities
    \item \textbf{Examination:} body mass index, neck circumference, blood pressure, and inspection of the nose, oropharynx, and jaw position
    \item \textbf{Home sleep study:} a portable monitor recording nasal airflow, oxygen saturation, heart rate, and respiratory effort during a night at home
    \item \textbf{Polysomnography:} in-laboratory monitoring of sleep stages, breathing, oxygen saturation, and limb movements when the diagnosis is uncertain or central sleep apnoea is suspected
    \item \textbf{Blood tests:} thyroid function, fasting glucose or glycated haemoglobin, and lipid profile to screen for contributing conditions
    \item \textbf{Referral:} to a sleep physician or respiratory physician for interpretation and treatment planning in moderate or severe disease
\end{itemize}

\section*{Management}
Treatment aims to keep the airway open during sleep, relieve symptoms, and reduce the associated cardiovascular risk. The choice of therapy depends on the severity of the apnoea, the anatomy of the airway, and patient preference.

\begin{itemize}
    \item \textbf{Weight loss:} even a 5--10\% reduction in body weight can substantially reduce or resolve apnoea severity and should be actively supported
    \item \textbf{Continuous positive airway pressure (CPAP):} the mainstay treatment for moderate to severe OSA, delivered through a nasal or facial mask overnight
    \item \textbf{Mandibular advancement splint:} a custom-fitted dental device that holds the lower jaw and tongue forward, suitable for mild to moderate cases
    \item \textbf{Positional therapy:} side-sleeping strategies for people whose apnoeas occur mainly when supine
    \item \textbf{Lifestyle measures:} avoiding alcohol and sedatives in the hours before bed, quitting smoking, and treating nasal congestion
    \item \textbf{Children:} adenotonsillectomy is usually curative when tonsils and adenoids are enlarged
    \item \textbf{Central sleep apnoea:} may respond to treatment of the underlying cause, such as optimising heart failure therapy or reducing opioid use
    \item \textbf{Follow-up:} regular review to check adherence, adjust pressure settings or the splint, and monitor symptoms and cardiovascular risk
\end{itemize}

\section*{Complications}
Untreated obstructive sleep apnoea places sustained strain on the cardiovascular system and affects day-to-day safety and quality of life.

\begin{itemize}
    \item Excessive daytime sleepiness increases the risk of motor vehicle and workplace accidents
    \item Hypertension is common and often difficult to control without treating the apnoea
    \item Increased risk of ischaemic heart disease, stroke, heart failure, and atrial fibrillation
    \item Higher risk of type 2 diabetes and features of the metabolic syndrome
    \item Mood disturbance, depression, and impaired memory and concentration
    \item Relationship strain caused by loud snoring and fragmented sleep for both partners
    \item Increased peri-operative risk, particularly with sedative anaesthetic agents, so anaesthetists should be informed of the diagnosis
\end{itemize}

\section*{Prognosis and Outcomes}
Sleep apnoea is a chronic condition, but it responds well to treatment. With consistent use of CPAP or a mandibular advancement splint, daytime sleepiness typically improves within weeks, and blood pressure and cardiovascular risk can fall. Substantial weight loss can markedly reduce severity and, in some cases, resolve the condition altogether. Children usually recover completely after removal of enlarged tonsils and adenoids. For most adults the underlying predisposition persists, so ongoing treatment and regular monitoring are required, but the majority manage the condition successfully and return to normal daytime alertness.

\section*{Prevention and Self-Care}
\begin{itemize}
    \item Maintain a healthy weight through a balanced diet and regular physical activity
    \item Avoid alcohol and sleeping tablets in the hours before bedtime
    \item Stop smoking and avoid exposure to second-hand smoke
    \item Sleep on your side rather than your back if your apnoea is position-dependent
    \item Treat allergies or nasal congestion that block the nose
    \item If you use CPAP, wear it every night and have the mask fitted correctly to maximise comfort and adherence
    \item Keep the sleep environment dark, quiet, and cool, and maintain a regular sleep-wake schedule
    \item For children, seek early review for persistent snoring, mouth breathing, or restless sleep
\end{itemize}

\section*{Sources and Further Reading}
The following organisations provide reliable, evidence-based information on sleep apnoea for patients and healthcare students.

\begin{itemize}
    \item NHS. \textit{Sleep apnoea: symptoms, diagnosis, and treatment.} https://www.nhs.uk/conditions/sleep-apnoea/
    \item Mayo Clinic. \textit{Obstructive sleep apnoea: overview and treatment options.} https://www.mayoclinic.org/diseases-conditions/obstructive-sleep-apnea/
    \item MSD Manual. \textit{Obstructive sleep apnoea (OSA).} https://www.msdmanuals.com/professional/pulmonary-disorders/sleep-apnea/obstructive-sleep-apnea
    \item MedlinePlus. \textit{Sleep apnoea.} https://medlineplus.gov/sleepapnea.html
    \item National Heart, Lung, and Blood Institute (NHLBI). \textit{Sleep apnoea information.} https://www.nhlbi.nih.gov/health/sleep-apnea
    \item World Health Organization. \textit{Sleep and sleep disorders.} https://www.who.int/health-topics/sleep
\end{itemize}
